\begin{figure}[ht]
\centering
\begin{tikzpicture}[x=1cm, y=1cm, font=\small]

% Parabola y = 0.35 x^2 on x in [-2.6, 2.6], with osculating circles at
% the vertex (large radius implied by small curvature) and at a flank
% (smaller radius). Plot scale: x by 1.4, y by 1.4.

% Axes
\draw[->, thick] (-4.2, 0) -- (4.4, 0) node[right] {$x$};
\draw[->, thick] (0, -0.4) -- (0, 4.6) node[above] {$y$};

% Parabola: plotx = 1.4*x, ploty = 1.4*(0.35*x*x)
\draw[very thick, engblue] plot[domain=-2.6:2.6, samples=80]
  ({1.4*\x}, {1.4*0.35*\x*\x});

% Osculating circle at the vertex (x=0). For y = a x^2, kappa(0) = 2a,
% radius R = 1/(2a) = 1/0.7 approx 1.43 in data units; centre at (0, R).
% Scaled: centre (0, 1.4*1.43) = (0, 2.0), radius 1.4*1.43 = 2.0.
\draw[thick, engred, dashed] (0, 2.0) circle (2.0);
\fill[engred] (0, 2.0) circle (1.4pt);
\node[engred, anchor=west, font=\scriptsize] at (0.15, 2.0)
  {centre, vertex};
\fill[engblack] (0, 0) circle (1.6pt);
\node[engblack, anchor=north east, font=\scriptsize] at (-0.1, -0.05)
  {$R_{0} = 1/(2a)$};

% Osculating circle at a flank point x = 1.8. y = 0.35*1.8^2 = 1.134.
% y' = 0.7*1.8 = 1.26; y'' = 0.7. kappa = |y''|/(1+y'^2)^{3/2}
%   = 0.7/(1+1.5876)^{1.5} = 0.7/4.158 = 0.1684; R = 5.94 data units.
% This is large; to keep the figure compact we draw a representative
% arc near the point rather than the full circle.
% Tangent direction at x=1.8: slope 1.26, unit normal points up-left.
% Normal direction (data): (-1.26, 1)/|.| = (-0.783, 0.622).
% Centre (data) = (1.8, 1.134) + R*(-0.783, 0.622)
%   = (1.8 - 5.94*0.783, 1.134 + 5.94*0.622) = (-2.85, 4.83).
% Scaled centre: (1.4*-2.85, 1.4*4.83) = (-3.99, 6.76); radius 1.4*5.94=8.3.
% Draw only a short arc passing through the scaled flank point
% (1.4*1.8, 1.4*1.134) = (2.52, 1.588).
\fill[engblack] (2.52, 1.588) circle (1.6pt);
\node[engblack, anchor=west, font=\scriptsize] at (2.6, 1.45)
  {flank point};
\draw[thick, engorange, dashed]
  ([shift={(-3.99,6.76)}] 230:8.3) arc (230:255:8.3);
\node[engorange, anchor=south west, font=\scriptsize] at (2.7, 1.9)
  {$R_{1} > R_{0}$};

% Tangent at flank point
\draw[thin, enggray] (2.52 - 1.2*1.4, 1.588 - 1.2*1.26)
  -- (2.52 + 0.7*1.4, 1.588 + 0.7*1.26);

\node[engblue, anchor=west, font=\scriptsize] at (-4.0, 4.3)
  {$y = a x^{2}$};

\end{tikzpicture}
\caption[Radius of curvature and the osculating circle]{The
osculating circle at a point on $y = a x^{2}$ shares the curve's
position, tangent, and second derivative there. At the vertex the
curvature $\kappa = 2a$ is largest and the radius $R_{0} = 1/(2a)$
is smallest; the circle hugs the curve tightly. At the flank point
the slope $y'$ is no longer zero, the denominator
$(1 + (y')^{2})^{3/2}$ in the curvature formula grows, the curvature
drops, and the matching radius $R_{1} > R_{0}$ widens. A vehicle
tracing the curve at constant speed feels a lateral acceleration
$v^{2}/R$ that peaks where the radius is smallest.}
\label{fig:vol02:ch05:curvature}
\end{figure}
